% HW: PAA 1-4
% CW: PAA 5-7, 9

\documentclass{article}

\usepackage[margin = 1in]{geometry}

\usepackage{graphicx, amsmath}
\graphicspath{ {./images/} }

\title{Ch. 11 Solutions}
\author{Principles of Microeconomics}
\date{\today}

\begin{document}

\maketitle

\section*{Quick Quiz}

\begin{enumerate}

\item (a) -- By definition, private goods are excludable and rival, and club goods are excludable and non-rival. 

\item (b) -- By definition, private goods are excludable and rival, and common resources are non-excludable and rival.

\item (b) -- National defense is non-excludable because once national defense is provided it is impossible to prevent anyone who didn't pay for it from benefitting from the protection. It's non-rival because one person benefitting from the protection doesn't diminish anyone else's ability to benefit from the protection.

\item (b) -- Because they're non-excludable, people can benefit from them without paying for them. That incentivizes people to try and free ride which makes it difficult to provide public goods profitably. As a result, the private market tends not to provide public goods.

\item (d) -- The total value to residents is
	\begin{equation*}
	\$80 + \$50 - \$30 = \$100.
	\end{equation*}
	That's less than the total cost, so it's more efficient not to provide the fireworks display. 
	
\item (d) -- Fish in the ocean are non-excludable because it's difficult to prevent people from catching them. They are rival because if one person catches a fish, there are fewer fish for other people to catch. 

\item (c) -- The fact that one person's use of a common resource diminishes another person's ability to use the common resource is a negative externality. The person using the good has no incentive to account for their negative impact on someone else, so they are likely to overuse the good.

\item (d) -- During other times of day, the road isn't congested, so one person using the road doesn't slow down everybody else on the road. Since it's non-rival, it isn't at risk of overuse, so there's no reason to disincentivize people from using it.

\end{enumerate}

\section*{Questions for Review}

\begin{enumerate}

\item A good is excludable if someone can be prevented from consuming it. A good is rival in consumption if one person's consumption of a unit diminishes another person's consumption of that unit.

A slice of pizza is excludable because a chef can refuse to give a potential customer any pizza. It's also rival in consumption because when one person consumes a unit of pizza, another person can't consume that same unit of pizza.   

\item A public good is a non-excludable, non-rival good.National defense is an example. Once one person pays for national defense, they can't prevent anyone else from using the protection, so it's non-excludable. One person using the protection doesn't diminish anybody else's use of the protections, so it's also non-rival.

Private markets generally can't provide public goods. If one person were to buy a public good, everyone else would benefit without paying for it. Because people can benefit without paying, people try and free ride. As a result, it's difficult to provide a public good profitably, so private markets generally don't. 

\item The cost-benefit analysis of a public good is the attempt to compare the total cost and total benefit of providing a public good.

It's important because it's the best way of deciding which public goods to provide and how much of them to provide.

It's difficult because you have to estimate how much everyone will benefit from the good. Unlike with private goods, people don't have to pay for public goods, so we can't estimate people's benefit based on their willingness to pay. Additionally, we can't ask them how much they benefit because they have an incentive to lie based on whether they want the public good.

\item A common resource is a non-excludable, rival good. Seats on the subway are an example. They're non-excludable because you can't prevent anyone from using them. They're rival because when one person takes a seat, another person can't take the same seat. 

Without government intervention people will use common resources too much because they have a negative externality. When one person uses a common resource and diminishes another person's ability to use it, the other person is made worse off without compensation. That external cost of using a common resource is a negative externality, so the market quantity of a common resource will be greater than the socially optimal quantity.

\end{enumerate}

\section*{Problems and Applications}

\begin{enumerate}

\item

	\begin{enumerate}
	
	\item
	
		\begin{itemize}
		
		\item Police protection is a club good because it’s excludable but not rival. You can deny police protection to certain people, but one person’s use of police protection doesn’t usually diminish another person’s. 
		
		\item Snow plowing is a public good because it’s not excludable or rival. You can’t prevent people from using a plowed street, and one person’s use of the plowed street doesn’t usually diminish another person’s.
		
		\item Education is a private good because it’s excludable and rival. You can prevent people who don’t pay for an education from receiving it, and since classrooms can only fit so many kids, one person’s consumption of education prevents another person from consuming it.
		
		\item Rural roads are public goods because they’re not excludable or rival. You can’t prevent people from using a rural road, and there’s no traffic in rural areas, so one person’s use doesn’t diminish another person’s
		
		\item City streets are common resources because they’re rival but not excludable. They’re not excludable because you can’t prevent people from using them, but they’re rival because there’s so much traffic in cities that one person’s use diminishes another person’s by adding to congestion.
		
		\end{itemize}
	
	\item If the private sector doesn’t produce enough of a good because it has a positive externality, the government can increase production by producing the good itself. This is the case with education and police protection. Sometimes the public uses a government provided public good so much that it becomes a common resource. This is the case with city streets.
	
	\end{enumerate}
	
\item 

	\begin{enumerate}
	
	\item Public goods typically have positive externalities, so the free-market quantity is less than the efficient quantity. For example, national defense provides protection for all of the people inside a country — not just the people who pay for it. As a result, it has a positive externality.
	
	\item Common resources typically have negative externalities, so the free-market quantity is greater than the efficient quantity. For example, when someone uses a city street during rush hour, they create a negative externality because they create congestion that slows other people down.
	
	\end{enumerate}
	
\item

	\begin{enumerate}
	
	\item Fredo is a free rider.
	
	\item The government can provide non-excludable goods and pay for them by increasing taxes.
	
	\item The private market can provide TV in an excludable way like through a cable, satellite, or streaming service. That way people have to pay for using the service. Any of these solutions turn TV from a public good into a club good.
	
	\end{enumerate}
	
\item

	\begin{enumerate}
	
	\item It’s a public good because it’s not excludable or rival. It’s free, so everyone can use it without paying, and it doesn’t have many users so no one person’s use slows down another person’s use.
	
	\item It’s a common resource because it’s become rival. One person’s use diminishes another’s.
	
	\item People might overuse the internet because they don’t account for the costs they create for others when they decide how much to use it. One way to correct that is by adding a paywall and thereby turning it into a private good.
	
	\end{enumerate}

\item

	\begin{enumerate}

	\item It is. Once a movie is on the TV, it's difficult to prevent people in the room who didn't pay for the movie from watching, so a movie showing is not excludable. Additionally, one person watching the movie doesn't diminish another person's ability to watch the movie, so a movie showing is not rival. 

	\item They should stream three movies. Their collective willingness to pay drops below \$8 after the third movie.
	
	\item They would spend $\$8 \cdot 3 = \$24$ total, so each person would spend \$6. Therefore, each person's consumer surplus would be:
	
		\begin{align*}
		\text{Dwayne: } &CS = 7 + 6 + 5 - 6 = 12 \\
		\text{Javier: } &CS = 5 + 4 + 3 -6 = 6 \\
		\text{Salman: } &CS = 3 + 2 + 1 - 6 = 0 \\
		\text{Chris: } &CS = 2 + 1 - 6 = -3
		\end{align*}
		
	\item Everyone could pay the same fraction of their willingness to pay The difficulty is getting people to tell the truth about their willingness to pay if they know it will determine what share of the cost they'll bear.
	
	\item He will not. He'll have an incentive to overstate his willingness to pay in an effort to get them to watch more movies. Each additional movie only costs him \$2 because the roommates split the \$8 cost four ways. His willingness to pay for a movie is greater than \$2 for a fourth and fifth movie, so he wants them watch at least five movies.
	
	\item It's difficult to find the optimal number of movies and then pay for them in a fair way.

	\end{enumerate}
	
\item

	\begin{enumerate}
	
	\item Basic research is a public good because it's not excludable or rival. Once it's published, you can prevent people who didn't fund the research from using it. Additionally, one person's use of the research doesn't diminish another person's ability to use it.
	
	\item The United States funds a lot of research at universities which it pays for using tax revenue.
	
	\item It's not totally consistent with my argument that basic research is not excludable. You could argue that there are barriers to foreign firms using domestic research. For one, if the research is in English, it will take time to translate the findings into other languages. For two, firms that are geographically closer to the researchers might have more exposure to their findings and an easier time adopting the advances. As a result, basic research might be slightly excludable. That said, over time the research will desseminate, and it will become non-excludable. 
	
	\end{enumerate}
	
\item

	\begin{enumerate}
	
	\item They would. The collective willingness to pay for the fireworks exceeds the cost.
	
	\item If they split the cost equally, everyone would pay \$120 for the fireworks. In that case, Callie would vote in favor because her willingness to pay is greater than \$120, but everyone else would vote no because their willingness pay is less than \$120. As a result, the vote would not yield the same answer as the cost-benefit analysis.
	
	\item No, the collective willingness to pay for the fireworks is less than the cost.
	
	\item Again, each person would pay \$120. Ned and Bess would both vote in favor because their willingness to pay is greater than \$120. Nancy would vote against because her willingness to pay is less than \$120. As a result, the vote would not yield the same answer as the cost-benefit analysis.
	
	\item It's difficult to decide whether or not to supply them. 
	
	\end{enumerate}

\setcounter{enumi}{8}	
\item During rush hour, the transportation system becomes overcrowded, so not everyone can fit on the trains. That means the trains become rival, and one person's use diminishes another person's ability to use them. In other words, the trains become a private good that get overused and create a negative externality. By increasing the price of using the transportation system during rush hour, the government is reducing the free-market quantity of train usage down towards its efficient quantity.

\end{enumerate}

\end{document}